\begin{corollary}{3.11.3}\label{II.3.11.3}
Let \(u\colon X\to Y\) be an \(S\)-morphism. One has an isomorphism
functorial in~\(\mathcal{M}\) (where, in the right-hand term, \(X\) lies over
\(Y\) via~\(u\)):
\[
L^{u}_{\Hom_{S}(X,Y)/S}(\mathcal{M})
  \isomto
\Hom_{Y/S}\bigl(X,\,T_{Y/S}(\mathcal{M})\bigr).
\]
{\renewcommand{\thefootnote}{(38)}\nde{The following sentence has been added.}}
It is an isomorphism of \(\mathbf{O}_{S}\)-modules if \(Y/S\) satisfies~(E).
\end{corollary}

In particular, for \(Y=X\), \(\End_{S}(X)\) is an \(S\)-functor in monoids, hence
\emph{a~fortiori} an \(S\)-H-functor; recalling that
\(\Lie(\End_{S}(X)/S,\mathcal{M})\) denotes
\(L^{e}_{\End_{S}(X)/S}(\mathcal{M})\), where \(e\) is the unit section
(\cf 3.9.0.1), one obtains:

\begin{corollary}{3.11.4}\label{II.3.11.4}
One has an isomorphism functorial in~\(\mathcal{M}\)
\[
\Lie(\End_{S}(X)/S,\mathcal{M})
  \isomto
\prod_{X/S}T_{X/S}(\mathcal{M})\,;
\]
{\renewcommand{\thefootnote}{(38)}\footnotemark[38]}
it is an isomorphism of \(\mathbf{O}_{S}\)-modules if \(Y/S\) satisfies~(E).
\end{corollary}

\begin{remark}{3.11.5}\label{II.3.11.5}
{\renewcommand{\thefootnote}{(39)}\nde{This remark has been added.}}
Suppose that \(X/S\) satisfies~(E). Then
\[
\prod_{X/S}T_{X/S}(\mathcal{M})
  =\Hom_{X/S}\bigl(X,\,T_{X/S}(\mathcal{M})\bigr)
\]
is endowed with a structure of
\(\bigl(\prod_{X/S}\mathbf{O}_{X}\bigr)\)-module, \ie for every \(S'\to S\),
\begin{equation*}
\begin{aligned}
&\Hom_{X/S}\bigl(X,\,T_{X/S}(\mathcal{M})\bigr)(S') \\
&\qquad
=\bigl\{\psi\in\Hom_{X}\bigl(I_{S'}(\mathcal{M})\times_{S}X,\,X\bigr)
  \bigm|
  \psi\circ(\varepsilon_{\mathcal{M}}\times\mathrm{id}_{X})
  =\mathrm{pr}_{X}\bigr\}
\end{aligned}
\end{equation*}
is endowed with a structure of \(\mathbf{O}(X\times_{S}S')\)-module, functorial
in~\(S'\). This follows, at one's choice, from~3.6 and from the properties of
the functor \(\prod_{X/S}\) (\cf 1.2), or else from the proof of~3.11.1.
\end{remark}

We shall now interpret geometrically the definition of the tangent bundle.
{\renewcommand{\thefootnote}{(40)}\nde{The original has been simplified in what follows.}}
Let \(U\) be an \(S\)-functor; by~I 1.7.2, one has isomorphisms functorial
in~\(\mathcal{M}\)
\begin{align*}
T_{X/S}(\mathcal{M})(U)
  &=\Hom_{S}\bigl(U,\,\Hom_{S}(I_{S}(\mathcal{M}),X)\bigr) \\
  &\simeq\Hom_{S}\bigl(I_{S}(\mathcal{M}),\,\Hom_{S}(U,X)\bigr) \\
  &\simeq\Hom_{I_{S}(\mathcal{M})}
      \bigl(U_{I_{S}(\mathcal{M})},\,X_{I_{S}(\mathcal{M})}\bigr).
\end{align*}
In particular, the morphism \(\mathcal{M}\to 0\) gives a commutative diagram:
\[
\xymatrix{
\Hom_{S}\bigl(U,\,T_{X/S}(\mathcal{M})\bigr)
  \ar[r]^{\sim}
  \ar[d]_{\pi_{\mathcal{M}}}
&
\Hom_{I_{S}(\mathcal{M})}
  \bigl(U_{I_{S}(\mathcal{M})},\,X_{I_{S}(\mathcal{M})}\bigr)
  \ar[d]
\\
\Hom_{S}(U,X)
  \ar[r]^{\textup{identity}}
&
\Hom_{S}(U,X)\,,
}
\]
where the second vertical arrow is obtained by the base change
\(\varepsilon_{\mathcal{M}}\colon S\to I_{S}(\mathcal{M})\).
{\renewcommand{\thefootnote}{(41)}\nde{Via the isomorphism
\(\Hom_{I_{S}(\mathcal{M})}
  \bigl(U_{I_{S}(\mathcal{M})},\,X_{I_{S}(\mathcal{M})}\bigr)
  \simeq
  \Hom_{S}\bigl(U\times_{S}I_{S}(\mathcal{M}),\,X\bigr)\),
this is equivalent to Remark~3.1.1.
Likewise, 3.12 is equivalent to the first part of Remark~3.4.2.}}

Consequently:

\begin{proposition}{3.12}\label{II.3.12}
Let \(h_{0}\colon U\to X\) be an \(S\)-morphism. Then
\(\Hom_{X}\bigl(U,\,T_{X/S}(\mathcal{M})\bigr)\) is identified with the set of
\(I_{S}(\mathcal{M})\)-morphisms from \(U_{I_{S}(\mathcal{M})}\) into
\(X_{I_{S}(\mathcal{M})}\) which restrict to~\(h_{0}\) on~\(U\) (viewed as a
subobject of \(U\times_{S}I_{S}(\mathcal{M})\) via
\(\mathrm{id}_{U}\times_{S}\varepsilon_{\mathcal{M}}\)).
\end{proposition}

{\renewcommand{\thefootnote}{(42)}\nde{The original has been spelled out in what follows.}}
In particular, for \(U=X\) and \(h_{0}=\mathrm{id}_{X}\), one obtains:

\begin{corollary}{3.12.1}\label{II.3.12.1}
The set \(\Gamma\bigl(T_{X/S}(\mathcal{M})/X\bigr)\) is identified with the set
of \(I_{S}(\mathcal{M})\)-endomorphisms \(\varphi\) of \(X_{I_{S}(\mathcal{M})}\)
which induce the identity on~\(X\), \ie such that the diagram below is
commutative:%
{\renewcommand{\thefootnote}{(43)}\nde{When \(\mathcal{M}=\OX\), these are
called ``infinitesimal endomorphisms'' of~\(X\); see~6.2 at the end of this
expos\'e.}}
\[
\xymatrix{
I_{X}(\mathcal{M}) \ar[r]^{\varphi}
  & I_{X}(\mathcal{M}) \\
X \ar[u]^{\varepsilon_{\mathcal{M}}} \ar[r]^{\mathrm{id}}
  & X \ar[u]^{\varepsilon_{\mathcal{M}}}
}
\]
\end{corollary}

{\renewcommand{\thefootnote}{(44)}\nde{The original has been spelled out in what follows.}}
On the other hand, by~3.11.2,%
% typo?: kept as printed: ``3.11.2'' (the isomorphism is the content of 3.11.4)
\(\Gamma\bigl(T_{X/S}(\mathcal{M})/X\bigr)
  \simeq
  \Lie(\End_{S}(X)/S,\mathcal{M})(S)\);
if moreover \(X/S\) satisfies~(E), then \(\End_{S}(X)/S\) satisfies~(E), so
\(\Lie(\End_{S}(X)/S,\mathcal{M})\) is an \(\mathbf{O}_{S}\)-module, by~3.6
(and in fact a \(\bigl(\prod_{X/S}\mathbf{O}_{X}\bigr)\)-module, by~3.11.5).
Applying~3.9, one deduces from this:

\begin{proposition}{3.13}\label{II.3.13}
If \(X/S\) satisfies~(E), the abelian group
\(\Gamma\bigl(T_{X/S}(\mathcal{M})/X\bigr)\) is identified with the set of
\(I_{S}(\mathcal{M})\)-endomorphisms of \(X_{I_{S}(\mathcal{M})}\) which induce
the identity on~\(X\). Consequently, every \(I_{S}(\mathcal{M})\)-endomorphism
of \(X_{I_{S}(\mathcal{M})}\) which induces the identity on~\(X\) is an
automorphism.
\end{proposition}

\begin{corollary}{3.13.1}\label{II.3.13.1}
Let \(u\colon X\to Y\) be an \(S\)-isomorphism, with \(Y/S\) satisfying~(E).
Every \(I_{S}(\mathcal{M})\)-morphism from \(X_{I_{S}(\mathcal{M})}\) into
\(Y_{I_{S}(\mathcal{M})}\) which extends~\(u\) is an isomorphism.
\end{corollary}

\begin{corollary}{3.13.2}\label{II.3.13.2}
If \(Y/S\) satisfies~(E), then the monomorphism
\[
\Isom_{S}(X,Y)\to\Hom_{S}(X,Y)
\]
induces, for every
\(u\in\Isom_{S}(X,Y)\), an isomorphism
\[
L^{u}_{\Isom_{S}(X,Y)/S}(\mathcal{M})
  \isomto
L^{u}_{\Hom_{S}(X,Y)/S}(\mathcal{M}).
\]
\end{corollary}

\begin{proof}
{\renewcommand{\thefootnote}{(45)}\nde{The following proof has been added.}}
One must see that
\(L^{u}_{\Isom_{S}(X,Y)/S}(\mathcal{M})(S')
  \to
  L^{u}_{\Hom_{S}(X,Y)/S}(\mathcal{M})(S')\)
is a bijection, for every \(S'\to S\). By base change (\cf 3.4), it suffices
to do this for \(S'=S\). In this case,
\(L^{u}_{\Hom_{S}(X,Y)/S}(\mathcal{M})(S)\)
(\resp \(L^{u}_{\Isom_{S}(X,Y)/S}(\mathcal{M})(S)\))
is the set of \(I_{S}(\mathcal{M})\)-morphisms
(\resp automorphisms)
\(X_{I_{S}(\mathcal{M})}\to Y_{I_{S}(\mathcal{M})}\)
which extend~\(u\), and one applies the preceding corollary.
\end{proof}

\begin{corollary}{3.13.3}\label{II.3.13.3}
If \(X/S\) satisfies~(E), the monomorphism
\(\Aut_{S}(X)\to\End_{S}(X)\)
induces, for every \(u\in\Aut_{S}(X)\), an isomorphism
\(L^{u}_{\Aut_{S}(X)/S}(\mathcal{M})
  \isomto
  L^{u}_{\End_{S}(X)/S}(\mathcal{M})\).
In particular, one has
\[
\Lie(\Aut_{S}(X)/S,\mathcal{M})
  \isomto
\Lie(\End_{S}(X)/S,\mathcal{M})
  \isomto
\prod_{X/S}T_{X/S}(\mathcal{M})\,,
\]
{\renewcommand{\thefootnote}{(46)}\nde{The following has been added.}}
so that \(\Lie(\Aut_{S}(X)/S,\mathcal{M})\) is endowed with a structure of
\(\bigl(\prod_{X/S}\mathbf{O}_{X}\bigr)\)-module.
\end{corollary}

\numpara{3.14}\label{II.3.14}
Suppose, to conclude, that \(X\) is representable.%
{\renewcommand{\thefootnote}{(47)}\nde{The reminder of~2.2.2 has been added, and the sequel has been spelled out.}}
In this case, one has seen in~2.2.2 that the \(X\)-functor \(T_{X/S}\) is
representable by \(\mathbf{V}(\Omega^{1}_{X/S})\), whence the bijections:
\[
\Gamma(T_{X/S}/X)
  \simeq
\Hom_{X}(\Omega^{1}_{X/S},\,\OX)
  \simeq
\mathrm{D\acute{e}r}_{\OS}(\OX).
\]

This can also be deduced from what precedes, as follows. By~3.13,
\(\Gamma(T_{X/S}/X)\) is identified with the set of infinitesimal
endomorphisms of~\(X\) (\ie the \(I_{S}\)-endomorphisms of \(X_{I_{S}}\)
inducing the identity on~\(X\)). Now \(X\) and \(X_{I_{S}}\) have the same
underlying topological space, the corresponding sheaves of rings being
\(\OX\) and \(\mathcal{D}\OX=\OX\oplus\mathcal{M}\), where
\(\mathcal{M}=\OX\) is regarded as a square-zero ideal. Denoting by
\(\pi\colon\mathcal{D}\OX\to\OX\) the morphism of \(\OX\)-algebras which
vanishes on~\(\mathcal{M}\), one deduces from this that giving an
infinitesimal endomorphism of~\(X\) is equivalent to giving a morphism of
\(\OS\)-algebras \(\varphi\colon\OX\to\mathcal{D}\OX\) such that
\(\pi\circ\varphi=\mathrm{id}_{\OX}\), which is equivalent to giving an
\(\OS\)-derivation of the sheaf of rings~\(\OX\).

Moreover, one easily sees that if
\(D,D'\in\mathrm{D\acute{e}r}_{\OS}(\OX)\) and if one denotes
by~\(\varphi_{D}\) the infinitesimal endomorphism corresponding to~\(D\), then
\[
\varphi_{D+D'}=\varphi_{D}\circ\varphi_{D'}\,.
\]
This shows that the identification
\[
\{\text{infinitesimal endomorphisms of }X\}
  \simeq
\mathrm{D\acute{e}r}_{\OS}(\OX)
\]
is an isomorphism of abelian groups. Taking~3.13 (and~3.11.5) into account,
one has therefore constructed an isomorphism of abelian groups (and even of
\(\mathbf{O}(X)\)-modules)
\[
\Gamma(T_{X/S}/X)
  \isomto
\mathrm{D\acute{e}r}_{\OS}(\OX)
\]
which recovers the classical interpretation of tangent vector fields in terms
of derivations of the structure sheaf.%
{\renewcommand{\thefootnote}{(48)}\nde{Moreover, this isomorphism is functorial
in~\(S\): if \(S'\to S\), setting \(X'=X_{S'}\), one has
\(\Lie(\Aut_{S}(X)/S)(S')
  \simeq
  \Gamma(T_{X'/S'}/X')
  \simeq
  \mathrm{D\acute{e}r}_{\mathcal{O}_{S'}}(\mathcal{O}_{X'})\).}}
Let us moreover observe that \(\Gamma(T_{X/S}/X)\) is equal to
\(H^{0}(X,\mathfrak{g}_{X/S})\), where \(\mathfrak{g}_{X/S}\) is the dual
of~\(\Omega^{1}_{X/S}\).

\sgasection{4}{Tangent space to a group -- Lie algebras}
\label{II.sec.4}

\numpara{4.1}\label{II.4.1}
Let \(G\) be a group functor over~\(S\). By~3.9.0.2,
\(T_{G/S}(\mathcal{M})\) and \(\Lie(G/S,\mathcal{M})\) are endowed with
structures of groups over~\(S\), and one has group morphisms
\[
\xymatrix{
\Lie(G/S,\mathcal{M}) \ar[r]^{i}
  & T_{G/S}(\mathcal{M}) \ar[r]^{p}
  & G \ar@/_/[l]_{s}
}\,;
\]
by definition \(i\) is an isomorphism of \(\Lie(G/S)(\mathcal{M})\) onto the
kernel of~\(p\), and \(s\) is a section of~\(p\). It then follows from~I 2.3.7
that this sequence of morphisms allows one to identify
\(T_{G/S}(\mathcal{M})\) with the semidirect product of~\(G\) by
\(\Lie(G/S,\mathcal{M})\).

\begin{definition}{4.1.A}\label{II.4.1.A}
{\renewcommand{\thefootnote}{(49)}\nde{The numbering 4.1.A and 4.1.B has been
added in order to set off these definitions.}}
The corresponding operation of~\(G\) on \(\Lie(G/S,\mathcal{M})\) is denoted
\[
\Ad\colon G\longrightarrow
  \Aut_{S\text{-gr.}}\bigl(\Lie(G/S,\mathcal{M})\bigr)
\]
and called the \emph{adjoint representation} (relative to~\(\mathcal{M}\))
of~\(G\); one therefore has by definition, for \(x\in G(S')\) and
\(X\in\Lie(G/S,\mathcal{M})(S')\):
\[
\Ad(x)X=i^{-1}\bigl(s(x)\,i(X)\,s(x)^{-1}\bigr).
\]
If \(G\) is commutative, then so is \(T_{G/S}(\mathcal{M})\), and
\(\Ad(x)X=X\).
\end{definition}

\begin{definition}{4.1.B}\label{II.4.1.B}
{\renewcommand{\thefootnote}{(49)}\footnotemark[49]}
If \(G\) and \(H\) are two group functors over~\(S\) and if
\(f\colon G\to H\) is a group morphism, one deduces from it by functoriality
a morphism of exact sequences compatible with the canonical sections:
\[
\xymatrix{
1 \ar[r]
  & \Lie(G/S,\mathcal{M}) \ar[r] \ar[d]_{L(f)}
  & T_{G/S}(\mathcal{M}) \ar[r] \ar[d]_{T(f)}
  & G \ar[r] \ar[d]_{f}
  & 1
\\
1 \ar[r]
  & \Lie(H/S,\mathcal{M}) \ar[r]
  & T_{H/S}(\mathcal{M}) \ar[r]
  & H \ar[r]
  & 1
}\,;
\]
\(L(f)\), which one will also denote \(\Lie(f)\), is the \emph{derived
morphism} of~\(f\).%
{\renewcommand{\thefootnote}{(50)}\nde{If \(f\) is a monomorphism, then so are
\(L(f)\) and \(T(f)\), \cf 3.7.}}
\end{definition}

\begin{remark}{4.1.C}\label{II.4.1.C}
{\renewcommand{\thefootnote}{(51)}\nde{This remark has been added.}}
If \(G/S\) and \(H/S\) satisfy~(E), then \(L(f)\) respects the structures of
\(\mathbf{O}_{S}\)-modules deduced from the ``functoriality in~\(\mathcal{M}\)''
(\cf 3.6).
\end{remark}

\begin{proposition}{4.1.1}\label{II.4.1.1}
Let \(g\in G(S)\). Then
\(\Ad(g)\colon\Lie(G/S,\mathcal{M})\to\Lie(G/S,\mathcal{M})\)
is the derived morphism of \(\Int(g)\colon G\to G\).
\end{proposition}

Indeed \(\Ad(g)X=i^{-1}\bigl(\Int(g)\,i(X)\bigr)\), which is none other than
\(L(\Int(g))(X)\) by the very definition of the derived morphism.

Suppose that \(G/S\) satisfies~(E). Then, by proposition~3.9, the group
structure of \(\Lie(G/S,\mathcal{M})\) defined as above is none other than
the structure induced by its structure of \(\mathbf{O}_{S}\)-module (defined
by means of~(E)). One then deduces from the preceding proposition and from
the functoriality of the operation of~\(\mathbf{O}_{S}\) (\cf 3.6) the
corollary:

\begin{corollary}{4.1.1.1}\label{II.4.1.1.1}
Suppose that \(G/S\) satisfies~(E). Then \(\Ad\) sends \(G\) into the subgroup
\[
\Aut_{\mathbf{O}_{S}\text{-mod.}}\bigl(\Lie(G/S,\mathcal{M})\bigr)
\]
of \(\Aut_{S\text{-gr.}}\bigl(\Lie(G/S,\mathcal{M})\bigr)\), \ie for every
\(g\in G(S')\), \(\Ad(g)\) respects the structure of
\(\mathbf{O}_{S'}\)-module of \(\Lie(G'/S',\mathcal{M})\) (where
\(G'=G\times_{S}S'\)). In other words, \(\Ad\) is a linear representation
(\cf I,~3.2) of~\(G\) in the \(\mathbf{O}_{S}\)-module
\(\Lie(G/S,\mathcal{M})\).
\end{corollary}

\begin{remarks}{4.1.1.2}\label{II.4.1.1.2}
a)~In order that \(G/S\) satisfy~(E), it is necessary and sufficient that,
for every pair \((\mathcal{M},\mathcal{N})\) of free \(\OS\)-modules of finite
type, the diagram
\[
\xymatrix{
& \Lie(G/S,\,\mathcal{M}\oplus\mathcal{N}) \ar[dl] \ar[dr] & \\
\Lie(G/S,\mathcal{M}) \ar[dr]
  && \Lie(G/S,\mathcal{N}) \ar[dl] \\
& \Lie(G/S,0)=S &,
}
\]
obtained by applying the functor \(\Lie(G/S,\,)\) to the diagram \((*)\)
of~2.1, be cartesian.

b)~Suppose that \(G/S\) satisfies~(E). Then the derived morphism of the group
law \(\pi\colon G\times_{S}G\to G\) is none other than the addition law
in~\(\Lie(G/S,\mathcal{M})\). (N.B.\ \(\pi\) is not a group morphism, but
\(\pi(e,e)=e\), so the derived morphism \(L(\pi)\) sends
\(T^{(e,e)}_{(G\times_{S}G)/S}(\mathcal{M})
  =\Lie(G/S,\mathcal{M})\times_{S}\Lie(G/S,\mathcal{M})\)
into \(\Lie(G/S,\mathcal{M})\), \cf 3.7 and~3.8.)
For every \(n\in\ZZ\), one shows likewise that if one denotes by
\(n_{G}\colon G\to G\) the morphism of \(S\)-functors defined by
\(g\mapsto g^{n}\), then the derived morphism \(L(n_{G})\) is multiplication
by~\(n\) on \(\Lie(G/S)\), \cf 3.9.4.
\end{remarks}

\numpara{4.1.2.0}\label{II.4.1.2.0}
{\renewcommand{\thefootnote}{(52)}\nde{The original has been spelled out in
what follows, in order to exhibit the action of~\(G\) on the functors
\(\Hom_{G/S}\bigl(G,\,T_{G/S}(\mathcal{M})\bigr)\) and
\(\Hom_{S}\bigl(G,\,\Lie(G/S,\mathcal{M})\bigr)\).}}
Let us now consider the \(S\)-functor
\(\Hom_{G/S}\bigl(G,\,T_{G/S}(\mathcal{M})\bigr)\); for every \(S'\to S\),
one has
\(T_{G/S}(\mathcal{M})_{S'}\simeq T_{G_{S'}/S'}(\mathcal{M})\) (\cf 3.4) and
therefore:
\[
\Hom_{G/S}\bigl(G,\,T_{G/S}(\mathcal{M})\bigr)(S')
  \simeq
\Hom_{G_{S'}}\bigl(G_{S'},\,T_{G_{S'}/S'}(\mathcal{M})\bigr)
  =\Gamma\bigl(T_{G_{S'}/S'}(\mathcal{M})/G_{S'}\bigr).
\]
Let us first note that one has an isomorphism, functorial in~\(S'\),
\begin{equation*}
(*)\qquad
\Hom_{S'}\bigl(G_{S'},\,\Lie(G_{S'}/S',\mathcal{M})\bigr)
  \isomto
\Gamma\bigl(T_{G_{S'}/S'}(\mathcal{M})/G_{S'}\bigr)
\end{equation*}
which associates to every
\(f\colon G_{S'}\to\Lie(G_{S'}/S',\mathcal{M})\) the section
\(s_{f}\colon G_{S'}\to T_{G_{S'}/S'}(\mathcal{M})\) such that, for every
\(S''\to S'\) and \(g\in G(S'')\):
\[
s_{f}(g)=i\bigl(f(g)\bigr)\,s(g).
\]

Let \(h\) be an automorphism of the functor \(G_{S'}\) over~\(S'\), not
necessarily respecting the group structure. To every section~\(\tau\) of
\(T_{G_{S'}/S'}(\mathcal{M})\), one may associate \(h(\tau)\) defined by
transport of structure: it is for instance the only section of
\(T_{G_{S'}/S'}(\mathcal{M})\) making the diagram
\[
\xymatrix{
G_{S'} \ar[r]^{\tau} \ar[d]^{h}
  & T_{G_{S'}/S'}(\mathcal{M}) \ar[d]^{T(h)} \\
G_{S'} \ar[r]^{h(\tau)}
  & T_{G_{S'}/S'}(\mathcal{M})
}
\]
commutative.

In particular, let us take for~\(h\) the right translation \(t_{x}\) by an
element \(x\) of \(G(S')\), \ie \(h(g)=t_{x}(g)=g\cdot x\), for every
\(g\in G(S'')\), \(S''\to S'\). Then one immediately has
\[
t_{x}(s_{f})=s_{t_{x}(f)}
\]
where \(t_{x}(f)\) denotes the morphism from \(G_{S'}\) into
\(\Lie(G_{S'}/S',\mathcal{M})\) defined by
\[
t_{x}(f)(g)=f(g\cdot x^{-1})
\]
% typo?: kept as printed: S'\to S' (the parallel quantifier below is S''\to S')
for every \(g\in G(S'')\), \(S'\to S'\).

It follows that if one lets \(G\) operate by right translations on
\[
\Hom_{G/S}\bigl(G,\,T_{G/S}(\mathcal{M})\bigr)
  \qquad\text{and}\qquad
\Hom_{S}\bigl(G,\,\Lie(G/S,\mathcal{M})\bigr)
\]
in the following way: for every \(S'\to S\), \(x\in G(S')\),
\(\sigma\in\Gamma\bigl(T_{G_{S'}/S'}(\mathcal{M})/G_{S'}\bigr)\) and
\(f\in\Hom_{S'}\bigl(G_{S'},\,\Lie(G_{S'}/S',\mathcal{M})\bigr)\),
\[
(\sigma\cdot x)(g)=\sigma(g\cdot x^{-1})\cdot s(x)
  \qquad\text{and}\qquad
(f\cdot x)(g)=f(g\cdot x^{-1}),
\]
for every \(g\in G(S'')\), \(S''\to S'\), then the isomorphism \((*)\) above
respects the operations of~\(G\).

In particular, via this isomorphism, the elements of
\(\Hom_{G/S}\bigl(G,\,T_{G/S}(\mathcal{M})\bigr)^{G}(S')\)
correspond to the \emph{constant} morphisms from \(G_{S'}\) into
\(\Lie(G_{S'}/S',\mathcal{M})\) (\ie factoring through the projection
\(G_{S'}\to S'\)) or else to the elements of
\(\Lie(G_{S'}/S',\mathcal{M})(S')=\Lie(G/S,\mathcal{M})(S')\).

\par\addvspace{0.55\baselineskip}
\noindent\textbf{Terminology.}\ ---
The elements of
\(\Hom_{G/S}\bigl(G,\,T_{G/S}(\mathcal{M})\bigr)^{G}(S')\)
will be called ``sections of \(T_{G_{S'}/S'}(\mathcal{M})\)
\emph{invariant under right translation}''.

One then obtains:

\begin{proposition}{4.1.2}\label{II.4.1.2}
{\renewcommand{\thefootnote}{*}\footnote{The statements 4.1.2, 4.1.3
and~4.1.4 are obtained more simply by observing that the automorphisms
of~\(G\) invariant under right translations are the left
translations.\textsuperscript{(53)}}}%
The map
\(\Lie(G/S,\mathcal{M})(S)\to\Gamma\bigl(T_{G/S}(\mathcal{M})/G\bigr)\)
which associates to \(X\in\Lie(G/S,\mathcal{M})(S)\) the section
\(x\mapsto X\cdot x\) is a bijection of \(\Lie(G/S,\mathcal{M})(S)\) onto
the part of \(\Gamma\bigl(T_{G/S}(\mathcal{M})/G\bigr)\) consisting of the
sections invariant under right translation.
\end{proposition}
{\renewcommand{\thefootnote}{(53)}\ndetext{See for instance the proof of
propositions~4.6 and~2.5 of \cite{DG70}, \S II.4.}}
